
\subsubsection{2.1 Domain Mixture Optimization for Language Models}

The role of training data composition in downstream model performance is well-studied for pretraining. DoReMi \citep{Xie2023DoReMi} optimizes domain weights via Group Distributionally Robust Optimization, showing that the proportion of text from different domains substantially affects downstream performance. Chameleon \citep{Xie2025Chameleon} extends leverage-score weighting to heterogeneous corpora in both pretraining and finetuning contexts. DomainPilot \citep{Zhang2026DomainPilot} applies domain-level loss monitoring to optimize SFT data mixture, reporting +3.8\% improvement on LiveCodeBench over naive mixing.

These methods optimize mixture proportions from a fixed multi-source pool and do not isolate the effect of training on a single source in exclusion from others. The controlled comparison ''what happens when the model trains only on source A versus only on source B?'' is not addressed.

\subsubsection{2.2 Data Selection and Source Effects in Code SFT}

Code SFT papers report training data compositions but do not ablate source identity. WizardCoder \citep{Luo2023WizardCoder} and OctoPack \citep{Muennighoff2023OctoPack} improve HumanEval pass@1 by changing data quality and instruction format, but source identity is confounded with quality and format throughout. DeepSeek-Coder \citep{Guo2024DeepSeekCoder} reports fixed training compositions. Lv et al. [2025] show saturation effects in code SFT data selection but use a curated multi-source dataset without source isolation. Chen et al. [2024] ablate atomic versus synthetic training types --- the closest prior work --- but not source dataset identity with a cross-benchmark transfer matrix.

\subsubsection{2.3 Distributional Alignment as a Predictor of Fine-Tuning Performance}

Ben-David et al. [2010] provide theoretical foundations for the relationship between domain divergence and transfer performance in supervised learning. Zhang et al. [2025] (GRAPE) show that selecting instruction-tuning data by distribution alignment outperforms using three times as much unselected data --- the most direct precedent for our mechanistic hypothesis. GRAPE operates on general instruction-following without code-specialized embeddings or permutation-based statistical testing.

Our contribution extends this line to code-specific SFT: we provide the first permutation-tested evidence that CodeBERT embedding cosine similarity between training source and test benchmark rank-orders HumanEval+ pass@1 across four source conditions ($\rho$=1.0, p=0.042, n=10,000 permutations).

